% Adición acumulativa. Insertar después de la suspensión esturmiana existente.
% Conservar íntegro el texto y la figura fig:holonomia-esturmiana.
\subsection{Connexion nonadique, holonomie et représentation monodromique}
\label{mono-ampl:conexion}

La connexion régionale déjà construite est maintenant étudiée au moyen de trois objets : le transport de l'état complet, la représentation de ses tours et la suspension de ses itinéraires. Les opérations d'APP--TRIT--TPK et le catalogue initial conservent les définitions précédentes. La distinction nécessaire concerne ici ce que chaque représentation retient d'une même histoire.

À l'indice de prolongement \(k\), une coordinatisation de l'état conserve
\[
 \widehat x_k=(B_k,C_k,p_k,c_k,\Sigma_k,W_k,g_k,m_k).
\]
Ici, \(B_k\) réunit les trois blocs régionaux, \(C_k\) est le cylindre,
\(p_k\) le préfixe, \(c_k\) le registre de retenue arithmétique, \(\Sigma_k\) la signature
de frontière et \(W_k\) l'information incidencielle. La phase est
\(g_k=1+[(k-1)]_9\) ; \(m_k\) compte les tours complets. La réduction
\(\beta_k=[m_k]_{12}\) conserve la position dodécaphasique, tandis que
l'entier \(m_k\) distingue ses tours successifs.

\begin{definition}[Transport d'un tour nonadique]
\label{mono-ampl:def-vuelta}
Soient \(\mathcal R_k\subseteq\widehat X_k\times\widehat X_{k+1}\)
les relations d'extension admissible du TPK. La réalisation de
l'holonomie au niveau \(k\) est la composition relationnelle
\[
 \Gamma_{9,k}=\mathcal R_{k+8}\circ\cdots\circ\mathcal R_k.
\]
Pour \(s\geq1\), son prolongement à \(s\) tours est
\[
 \Gamma_{9s,k}=
 \Gamma_{9,k+9(s-1)}\circ\cdots\circ\Gamma_{9,k}.
\]
\end{definition}

La formulation relationnelle conserve les continuations admissibles. Sur
une histoire individualisée, chaque flèche mène à l'état effectivement
prolongé ; le bloc visible isolé peut encore admettre plusieurs
continuations. Cette différence est précisément la fonction de la mémoire et
des conditions de frontière dans la détermination de la trajectoire.

\begin{proposition}[Itération du retour et mémoire accumulée]
\label{mono-ampl:vueltas}
Dans toute composition admissible de \(s\) tours,
\[
 g_{k+9s}=g_k,\qquad m_{k+9s}=m_k+s,\qquad
 \beta_{k+9s}=\beta_k+s\pmod{12}.
\]
\end{proposition}
\begin{proof}
La règle d'un tour conserve la phase et incrémente le compteur d'une
unité. Son application à l'état déjà transporté répète exactement ces
deux opérations. La récurrence donne les deux premières égalités ; réduire la
seconde modulo \(12\) donne la troisième.
\end{proof}

Le compteur fournit ainsi un témoin entier de la différence entre les
tours. Après \(12\) tours, le couple des phases finies est rétabli, mais
le compteur augmente de \(12\). L'extension non scindée
\eqref{eq:extension108} exprime la composition du calendrier par son
cocycle ; la trajectoire conserve en outre le relèvement entier de cette donnée.
La conservation signifie ici le transport des contributions antérieures,
et non l'immobilité de toutes les coordonnées.

\subsubsection{Représentation bilatérale de l'indice de retour}

Le décalage \(T\) de \eqref{eq:weyl-relojes} admet une
expression en coordonnées de phase et de tour. Numérotons les phases de \(0\) à
\(8\), par un décalage de l'origine par rapport à \(g_k\), et
écrivons \(k=9m+r\), avec \(0\leq r<9\), sur la droite entière
qui revêt le cycle des phases. Sur
\(\mathcal H_{\rm ind}=\ell^2(\mathbb Z)\), la base
\(|k\rangle=|r,m\rangle\) exprime ce même décalage unitaire
\[
 T|r,m\rangle=
 \begin{cases}
 |r+1,m\rangle,&r<8,\\
 |0,m+1\rangle,&r=8.
 \end{cases}
\]
L'inverse réduit \(r\) d'une unité lorsque \(r>0\), et envoie
\(|0,m\rangle\) sur \(|8,m-1\rangle\). Le passage par l'origine de phase
incorpore donc l'unité exacte de quotient.

\Needspace{8\baselineskip}
\begin{proposition}[Représentation de la monodromie du cycle]
\label{mono-ampl:representacion}
Soit \(S=T^9\). Alors
\[
 S|r,m\rangle=|r,m+1\rangle,\qquad
 \mathbb Z\longrightarrow\mathcal U(\mathcal H_{\rm ind}),
 \quad s\longmapsto S^s
\]
est une représentation unitaire fidèle du groupe des tours du cycle.
\end{proposition}
\begin{proof}
Neuf décalages conservent le reste de la division par \(9\) et
incrémentent le quotient. Les puissances satisfont
\(S^{s+t}=S^sS^t\), et \(S^s|r,m\rangle=|r,m+s\rangle\) distingue
tout entier \(s\). La permutation d'une base orthonormée est unitaire.
\end{proof}

Cette représentation correspond au revêtement du cycle, dont le groupe
des tours est \(\mathbb Z\), et représente la phase avec son compteur. L'état
du TPK contient en outre les coordonnées de trajectoire déjà décrites.
L'extension bilatérale admet des indices négatifs comme coordonnées du
revêtement ; une trajectoire issue d'une condition initiale utilise
son demi-axe admissible. On dispose ainsi d'un inverse opératoriel
de l'indice sans transformer chaque relation d'extension en bijection de
l'état complet.

Avec \(D_M|k\rangle=\lfloor k/9\rfloor|k\rangle\), l'identité
\([D_M,S]=S\) exprime l'incrément de mémoire sur le domaine dense des
vecteurs à support fini. Si \(\mathcal I|k\rangle=|-k\rangle\) et
\(P_0\) projette sur les indices divisibles par \(9\), on obtient
également
\begin{equation}
 \mathcal I T\mathcal I=T^{-1},\qquad
 \mathcal I D_M\mathcal I=-D_M-(I-P_0).
 \label{mono-ampl:inversion-contador}
\end{equation}
En effet, pour \(k=9m+r\),
\(\lfloor-k/9\rfloor=-m\) si \(r=0\) et vaut \(-m-1\) si \(r>0\).
L'inversion conserve ainsi le terme de frontière de la division euclidienne.

\subsection{Horloge des blocs et horloge de capacité}
\label{mono-ampl:relojes}

L'indice d'une application de raffinement et le nombre de positions
décimales disponibles ont des lois liées, mais différentes. Pour
les exprimer, nous réservons \(n\) à l'indice des blocs ternaires et définissons
\[
 \rho=\log_{1000}729=\log_{10}9,\qquad
 \delta=1-\rho,\qquad
 \mathcal C_n=\lfloor(n+1)\rho\rfloor,\quad n\geq0.
\]
La lettre \(\mathcal C_n\) désigne ici la capacité, et non le registre
dodécaphasique \(K\). À l'origine de phase canonique, le mot mécanique
précédent fournit
\[
 Z_n=\sum_{j=0}^{n}z_j(0)=\lfloor(n+1)\delta\rfloor.
\]

\begin{proposition}[Bilan exact entre capacité et lacunes]
\label{mono-ampl:balance-capacidad}
Pour \(n\geq0\) et, dans la deuxième identité, \(n\geq1\),
\[
 \mathcal C_n=n-Z_n,\qquad
 \mathcal C_n-\mathcal C_{n-1}=1-z_n(0).
\]
\end{proposition}
\begin{proof}
L'irrationalité de \(\delta\) implique
\(\lceil(n+1)\delta\rceil=\lfloor(n+1)\delta\rfloor+1\).
Alors
\[
 \lfloor(n+1)(1-\delta)\rfloor
 =n+1-\lceil(n+1)\delta\rceil=n-Z_n.
\]
La soustraction de deux identités consécutives achève la preuve.
\end{proof}

Un événement \(z_n=1\) ajoute un bloc à l'histoire en maintenant la capacité
\(\mathcal C_n\) ; un événement \(z_n=0\) incrémente les deux compteurs. Sur
tout intervalle, les incréments de capacité et les lacunes ont pour somme
exactement sa longueur. En réduisant modulo \(q\), pour \(q=9\) ou
\(q=108\), on obtient
\[
 [\mathcal C_n]_q=[n]_q-[Z_n]_q.
\]
La phase uniforme des blocs et la phase de capacité sont liées par
la mémoire intégrée des lacunes. En particulier, une rétention de
capacité et un tour complet de la connexion ont des mécanismes
distincts, bien que tous deux puissent conserver une observation de phase.

L'inversion monotone de l'horloge donne un temps de première arrivée :
\[
 T_{\rm cap}(a)=\min\{n:\mathcal C_n=a\}
 =\left\lceil\frac a\rho\right\rceil-1,
 \qquad a\geq1.
\]
Avec \(\sigma=\rho^{-1}-1\), le temps nécessaire pour gagner \(b\) unités
de capacité est
\[
 T_{\rm cap}(a+b)-T_{\rm cap}(a)
 =b+\lceil(a+b)\sigma\rceil-\lceil a\sigma\rceil.
\]
La différence des plafonds prend l'un des deux entiers adjacents à
\(b\sigma\). On obtient ainsi \(9\) ou \(10\) blocs pour gagner
\(9\) unités de capacité, et \(113\) ou \(114\) pour en gagner \(108\).
Ce sont des temps de l'horloge inverse, mesurés depuis une première arrivée ; le
compteur des tours s'interprète toujours dans l'indice auquel il appartient.

\subsection{Suspension symbolique et cristal temporel apériodique}
\label{mono-ampl:cristal}

Le codage par rotation réunit périodicité de phase, récurrence locale
et apériodicité globale. Soit \(\mathcal X_\delta\subseteq\{0,1\}^{\mathbb Z}\)
l'adhérence, dans la topologie produit, des translations du mot
bilatéral \(z_n(\theta)\). Notons \(\sigma_{\rm sh}\) le
décalage de ses suites. En conservant une phase nonadique, le pas est
\[
 F(r,z)=([r+1]_9,\sigma_{\rm sh}z),\qquad
 (r,z)\in C_9\times\mathcal X_\delta.
\]

\begin{definition}[Cristal temporel apériodique symbolique]
\label{mono-ampl:def-cristal}
La suspension à toit unitaire du système précédent est l'espace
\[
 \mathcal S_\delta=
 \bigl((C_9\times\mathcal X_\delta)\times\mathbb R\bigr)/\sim,
 \qquad (y,t+1)\sim(Fy,t),
\]
avec le flot \(\Phi_s[y,t]=[y,t+s]\). On appelle \emph{cristal temporel
apériodique symbolique} ce modèle d'ordre temporel, accompagné de son
codage de lacunes et de la représentation du compteur des tours.
\end{definition}

L'adjectif temporel désigne l'action du flot et de sa section de retour ;
l'apériodicité caractérise ses itinéraires. La fréquence des uns dans chaque
suite est \(\delta\) : la borne uniforme de discrépance inférieure à un,
obtenue par télescopage dans toute fenêtre, persiste lors du passage à l'adhérence.
Une suite périodique aurait une fréquence rationnelle. Par conséquent,
\(\mathcal X_\delta\) ne contient pas d'orbites périodiques. La suspension
ne possède pas non plus d'orbites fermées : un retour du flot exigerait un temps
entier et une périodicité de l'itinéraire.

Les mots finis, en revanche, réapparaissent. Les itinéraires de longueur \(m\)
proviennent d'une partition de la circonférence par \(m+1\) points ; chaque
intervalle est visité de manière récurrente par la rotation irrationnelle. Même
lorsque les temps sont restreints à une phase fixée, le pas \(9\delta\) demeure
irrationnel. C'est pourquoi chaque mot apparaît dans les
\(9\) phases, et pourquoi valent les complexités déjà établies
\[
 p(m)=m+1,\qquad p_9(m)=9(m+1),\qquad p_3(m)=3(m+1).
\]
La récurrence des configurations locales coexiste avec l'absence d'une
période globale. Leur entropie topologique est nulle parce que ces complexités
croissent linéairement.

La réalisation par les observables \(V_9,W_\delta\) et le décalage
\(T\) exprime le même ordre par les covariances
\eqref{eq:weyl-relojes}. Les fréquences de l'horloge uniforme appartiennent à
\(\frac19\mathbb Z+\delta\mathbb Z\), modulo \(1\). L'horloge de
capacité conserve également la relation
\([\mathcal C_n]_9=[n]_9-[Z_n]_9\) ; elle n'est pas remplacée par la phase uniforme
lors de l'interprétation des retours. Le terme \emph{cristal temporel} est ainsi
défini au moyen d'un système dynamique mathématique. Sa réalisation physique
requiert la spécification d'une dynamique énergétique, de ses observables temporels
et de sa stabilité, données distinctes de la présente construction symbolique.

\subsection{Symétrie spéculaire des régions et double réalisation}
\label{mono-ampl:regiones}

L'échange spéculaire agit sur le bloc régional
\(B=(u^\pi,u^e,u^\varphi)\in(\mathbb F_3^6)^3\) par
\[
 \mathfrak cB=(u^e,u^\pi,u^\varphi).
\]
Son axe fixe et son agrégat sont respectivement \(u^\varphi\) et
\(s=u^\pi+u^e\). Cette fixation concerne l'involution dans une fibre ;
les deux coordonnées peuvent évoluer pendant le transport nonadique.
Avec \(q=u^\pi+u^e+u^\varphi\),
\(a=u^\pi+u^e-u^\varphi\) et \(d=u^\pi-u^e\), on retrouve
\[
 u^\varphi=2(q-a),\qquad s=q-u^\varphi,\qquad
 u^\pi=2(s+d),\qquad u^e=2(s-d).
\]
Les égalités sont composante par composante dans \(\mathbb F_3\), où
\(2^{-1}=2\). Le miroir conserve \(q,a,s\) et inverse \(d\) ; cette
coordonnée antisymétrique individualise la répartition que l'agrégat conserve
sans la distinguer. L'inversion \(\mathcal I\) de l'indice temporel et
l'involution \(\mathfrak c\) des canaux agissent donc sur des
données différentes.

Le retour de phase possède un témoin explicite :
\[
 \begin{aligned}
 B_1&=(012222\mid121221\mid112202),\\
 B_{10}&=(101222\mid112221\mid112002).
 \end{aligned}
\]
Les deux positions ont la phase \(1\) ; le bloc, l'axe et la mémoire ont
changé. De même, la synchronisation des dénombrements aux phases \(8\) et
\(9\) annule leurs différences relatives, tandis que la composante diagonale
passe de \((59,59,59)\) à \((43,43,43)\). L'égalité d'un observable
relatif n'épuise pas l'état qui le détermine.

La double réalisation conserve cette architecture. Dans la section
\ref{exc:doble-lectura}, le même préfixe détermine d'une part des cylindres
emboîtés et d'autre part une suite de mots codés avec leurs repères.
La première application passe à la limite par contraction des diamètres ; la seconde
satisfait \(r_{n,m}Q_n=Q_mp_{n,m}\) et définit \(Q_\infty\).
La figure \ref{mono-ampl:fig-regional} anticipe ces deux applications.
Le parcours incidenciel ultérieur construit des codes, des réseaux et la
réalisation de Moonshine dans leurs domaines propres ; le groupe des automorphismes
agit sur cette réalisation, et non sur les chiffres par une identification
implicite.

% Figura acumulativa; destino figures/monodromia_regional_ampliacion.tex.
\begin{figure}[!htbp]
\centering
\begin{tikzpicture}[x=1cm,y=1cm,>=Latex,font=\small,
  flow/.style={->,line width=.55pt},
  ann/.style={align=center,font=\footnotesize}]
\node[font=\bfseries] at (6.9,9.2)
  {Connexion nonadique et coordonnées régionales};
\begin{scope}[shift={(3.4,6.8)}]
 \foreach \k in {1,...,9}{
  \pgfmathsetmacro{\ang}{90-40*(\k-1)}
  \coordinate (p\k) at (\ang:1.52);
  \fill (p\k) circle (1.4pt);
  \node[fill=white,inner sep=1.1pt,font=\footnotesize]
    at (\ang:1.83) {\(\k\)};
 }
 \foreach \k in {1,...,8}{
  \pgfmathtruncatemacro{\next}{\k+1}
  \draw[flow,shorten >=3pt,shorten <=3pt] (p\k)--(p\next);
 }
 \draw[flow,shorten >=3pt,shorten <=3pt] (p9)--(p1);
 \node[ann] at (0,.1) {\(g_k\in C_9\)\\\(\Gamma_{9,k}\)};
 \node[ann] at (0,-2.23) {Un tour : \(g\mapsto g\), \(m\mapsto m+1\)};
\end{scope}
\node[ann,text width=5.25cm] at (10.1,7.75)
 {\(B_k=(u_k^\pi,u_k^e,u_k^\varphi)\)\\
 bloc régional dans une fibre};
\node at (8.55,6.5) {\(u^\pi\)};
\node at (11.65,6.5) {\(u^e\)};
\draw[<->,line width=.65pt] (8.95,6.5)--node[above]{\(\mathfrak c\)}(11.25,6.5);
\draw[densely dashed,gray] (10.1,5.72)--(10.1,7.23);
\node[fill=white,inner sep=2pt] at (10.1,5.6){\(u^\varphi\)};
\node[ann] at (10.1,4.6)
 {Axe \(u^\varphi\) et agrégat \(u^\pi+u^e\)\\
 invariants sous \(\mathfrak c\)};
\draw[flow] (5.45,6.8)--(7.25,6.8);
\draw[gray!60,line width=.3pt] (.55,4.02)--(13.25,4.02);
\node[ann,text width=12cm] (hist) at (6.9,3.47)
 {Histoire compatible : préfixes, cylindres, orientation, quotients,
 repères d'incidence et mémoire};
\coordinate (fork) at (6.9,2.85);
\draw[flow] (hist.south)--(fork);
\draw[flow] (fork)--(3.6,2.2);
\draw[flow] (fork)--(10.3,2.2);
\node[ann,text width=5.6cm] at (3.6,1.65)
 {Réalisation archimédienne\\
 \(I_{n+1}\subseteq I_n,\quad |I_n|=729^{-n}\)\\
 limite des préfixes régionaux};
\node[ann,text width=5.6cm] at (10.3,1.65)
 {Réalisation incidencielle\\
 \(r_{n,m}Q_n=Q_mp_{n,m}\)\\
 mots codés et repères compatibles};
\node[ann,text width=5.6cm] at (3.6,.42)
 {\(\pi,\varphi,e\) ; détermination dodécaphasique de \(\alpha\)};
\node[ann,text width=5.6cm] at (10.3,.42)
 {Witt--Golay--Leech--\(V^\natural\)\\
 automorphismes de la réalisation exceptionnelle};
\draw[flow] (3.6,1.03)--(3.6,.78);
\draw[flow] (10.3,1.03)--(10.3,.78);
\end{tikzpicture}
\caption{\fontsize{9}{11}\selectfont Neuf phases d'une connexion et deux réalisations de l'histoire
compatible. Le cycle représente la phase ; chaque tour incrémente le compteur.
Le diagramme régional représente l'involution \(\mathfrak c\), qui
échange clôture et propagation et fixe l'axe d'auto-échelle dans une fibre.
Cette invariance n'impose pas que l'axe soit constant pendant le transport.
Les deux applications inférieures utilisent le même préfixe : l'une détermine
sa limite numérique et l'autre conserve les mots et les repères d'incidence. Les
applications de cette seconde réalisation sont spécifiées dans la section
\ref{exc:doble-lectura} ; le parcours exceptionnel est développé ensuite.}
\label{mono-ampl:fig-regional}
\end{figure}


\clearpage
\subsubsection{Représentation rectiligne des phases et du retour}
\label{mono-recta:desarrollo}

La représentation sur le segment \([0,10]\) conserve l'origine
positionnelle d'APP et ordonne ses neuf marques intérieures. Dans ce diagramme,
les entiers \(1,\ldots,9\) sont des représentants de phase. La position
horizontale indique l'ordre des coordinatisations ; les étiquettes supérieures
spécifient les opérations régionales. La figure
\ref{mono-recta:figura} réunit ainsi la droite initiale, la fixation de
l'axe d'auto-échelle, la distribution spéculaire des deux canaux et le
retour avec mémoire.

% Recuperación de fig_nonadic_route_mini.tex de la síntesis académica.
% Los rótulos designan operaciones regionales; la suma se efectúa en F_3^6.
\begin{figure}[H]
\centering
\begin{tikzpicture}[x=1.12cm,y=1cm,font=\small,>=Latex,
  guide/.style={line width=.45pt,black!55}]
 \draw[line width=.8pt] (0,0)--(10,0);
 \foreach \k in {0,...,10}{
  \draw[guide] (\k,-.09)--(\k,.09);
  \node[below=2.5mm] at (\k,0) {$\k$};
 }
 \foreach \k in {1,...,9}{\fill (\k,0) circle (1.7pt);}
 \node[anchor=west,font=\footnotesize] at (0,3.05)
   {Segment générateur $[0,10]$ et neuf représentants de phase};
 \draw[guide] (5,.12)--(5,1.7);
 \node[above,align=center,font=\footnotesize] at (5,1.7)
   {fixation de l'axe\\$u^\varphi$};
 \draw[decorate,decoration={brace,amplitude=4pt},guide]
   (6,.65)--(8,.65);
 \node[above=5pt,align=center,font=\footnotesize] at (7,.65)
   {agrégat régional\\$u^e+u^\pi$};
 \draw[guide] (9,.12)--(9,1.7);
 \node[above,align=center,font=\footnotesize] at (9,1.7)
   {involution spéculaire\\$\pi\leftrightarrow e$};
 \draw[->,line width=.7pt] (9,-.7)
   .. controls (9,-1.55) and (1,-1.55) .. (1,-.7);
 \node[fill=white,inner sep=3pt,font=\footnotesize] at (5,-1.23)
   {retour $9\to1$;\quad $m\mapsto m+1$};
 \node[font=\footnotesize] at (5,-1.95)
   {$w_6\to w_{12}\to w_{18}\to w_{24}\to w_{30}\to R_{36}\to G_9$};
\end{tikzpicture}
\caption{Représentation rectiligne de la connexion nonadique. La marque
\(5\) situe la première saturation forte et l'axe régional d'autoéchelle ;
l'accolade \(6\)--\(8\) identifie l'agrégat ternaire
fixé dans chaque fibre ; la phase \(9\) sélectionne l'orientation de
clôture. Le retour conserve la phase et incrémente la mémoire. Les
abscisses numérotent les positions opératives, tandis que les étiquettes
\(u^\pi,u^e,u^\varphi\) désignent des blocs de \(\mathbb F_3^6\).}
\label{mono-recta:figura}
\end{figure}


La fixation de l'axe se formule dans la fibre d'une frontière donnée. Si
\(q=u^\pi+u^e+u^\varphi\) et
\(a=u^\pi+u^e-u^\varphi\), alors
\[
 u^\varphi=2(q-a),\qquad s=u^\pi+u^e=q-u^\varphi
 \quad\text{en }\mathbb F_3^6.
\]
Les données \((q,a)\) déterminent l'axe et l'agrégat ; la répartition ordonnée de \(s\) conserve la liberté spéculaire résolue par le prolongement. La cinquième phase possède une solution locale et une extension, avec la donnée de frontière \(c=111111\) et la signature résiduelle \(\rho_W=000\). C'est le sens de sa première saturation forte. Le symbole \(\varphi\) situé sur la marque \(5\) identifie l'axe régional qui intervient dans cette opération.

Les phases intermédiaires conservent cette décomposition au sein de chaque fibre,
tandis que leurs données de frontière évoluent. Le recensement complet des
multiplicités locales et des extensions de la branche distinguée est
\[
\begin{array}{c|rrrrrrrrr}
\text{phase}&1&2&3&4&5&6&7&8&9\\\hline
N_{\rm loc}&3&2&5&4&1&1&3&3&2\\
N_{\rm ext}&2&5&4&1&1&3&3&2&1
\end{array}
\]
La phase \(6\) combine unicité locale et trois prolongations ; la phase
\(7\) résout une fibre spéculaire triple ; la phase \(8\) synchronise
les trois recensements régionaux. La phase \(9\) conserve deux distributions
locales ayant le même axe et le même agrégat, dont l'une se
prolonge. Ces dénombrements expriment des fonctions différentes au sein d'une
même connexion et justifient les marques de la droite.

En particulier, pour les phases \(6,7,8\), les sommes régionales sont
respectivement \(012100\), \(101001\) et \(112000\). La conservation
représentée par l'accolade est une invariance par rapport à la répartition
spéculaire dans chaque fibre. L'évaluation réelle ultérieure applique les
lecteurs des cylindres compatibles aux histoires complètes ; la
somme interne représentée dans la figure appartient au domaine ternaire.

Enfin, l'arc \(9\to1\) représente le passage d'un tour au suivant. Les formules de la section \ref{mono-ampl:conexion} conservent le reste de phase et incrémentent le compteur de tours. Le diagramme rectiligne et le diagramme cyclique de la figure \ref{mono-ampl:fig-regional} expriment donc deux aspects complémentaires de la même holonomie : l'ordre des positions et la composition du retour sur l'état avec mémoire.


\subsection{Transport temporel et composition exponentielle tritique}
\label{mono-ampl:euler}

Les relations d'Euler de la section \ref{euler-trit-articulacion}
décrivent les réalisations locales des trois régimes. Pour une signature
fixée \(\tau\), le générateur satisfait \(J_\tau^2=-\tau I\) ;
la conjugaison algébrique envoie \(J_\tau\) sur \(-J_\tau\). Par conséquent,
\[
 \exp(sJ_\tau)\exp(tJ_\tau)=\exp((s+t)J_\tau),\qquad
 \overline{\exp(tJ_\tau)}=\exp(-tJ_\tau),
\]
et le produit d'une évolution par sa conjuguée est l'unité. Dans le type
elliptique, cette réalisation est donnée par \(\cos^2t+\sin^2t=1\) ; dans l'hyperbolique,
par \(\cosh^2t-\sinh^2t=1\) ; dans le parabolique,
par \((I+tN)(I-tN)=I\), puisque \(N^2=0\). La même loi de conservation
admet trois expressions quadratiques.

La réalisation elliptique ferme un tour ; la nilpotente conserve un
déplacement linéaire ; l'hyperbolique sépare expansion et contraction avec
un déterminant unitaire. Aux paramètres régionaux déjà générés,
\[
 \exp(\pi J_{\mathrm e})=-I,\qquad
 \exp(J_{\mathrm p})=I+J_{\mathrm p},\qquad
 \exp(2\log\varphi\,J_{\mathrm h})=F_{\rm av}^{\,2}.
\]
L'évaluation \(\exp(I)=eI\) reconnaît l'unité de la loi exponentielle
commune. Le nombre \(e\) n'est pas attribué exclusivement au régime parabolique.

Le transport d'une réalisation locale par un isomorphisme \(U\) conserve
ces relations : si \(J'=UJU^{-1}\), la série exponentielle donne
\(U\exp(tJ)=\exp(tJ')U\). L'identité transporte le régime et sa
norme. Une transition entre types quadratiques différents appartient, en
revanche, à la sélection des feuillets et des régimes du TPK ; ce n'est pas une conjugaison
qui modifie la valeur de \(J^2\). Le produit temporel conserve donc
l'ordre des transitions et l'information de leurs extrémités.

Cette articulation distingue la clôture d'une représentation locale de
l'holonomie d'une histoire. On peut avoir \(\exp(2\pi J_{\mathrm e})=I\)
en même temps que \(S|r,m\rangle=|r,m+1\rangle\). La première égalité
rétablit l'orientation du régime ; la seconde enregistre un autre tour. La
structure discrète conserve les deux : le retour local et la provenance
accumulée de l'état qui le réalise.
